\section{Zero correlation length}

\begin{definition}[Physical-trace transfer]
    \label{def:mpo_physical_trace_transfer}
    \lean{MPOTensor.physTraceTransfer,
        MPOTensor.verticalLoop_eq_physTraceTransfer}
    \leanok
    \uses{def:mpo_tensor}
    The physical-trace transfer of an MPO tensor $M$ is the virtual matrix
    obtained by contracting the ket and bra physical legs of one tensor:
    \begin{align}
        \mathcal T_M
         & =\sum_i M^{ii}.
        \notag
    \end{align}
    % Gap documented in docs/paper-gaps/cpsv16_zcl_canonical_form_normalization.tex.
    This is the transfer object appearing in the zero-correlation-length
    condition of \cite[Definition~4.2, lines~735--739]{Cirac2016MPDO_arXiv}.
\end{definition}

\begin{definition}[Literal physical-trace idempotence]
    \label{def:mpo_physical_trace_idempotence}
    \lean{MPOTensor.IsPhysicalTraceIdempotent}
    \leanok
    \uses{def:mpo_physical_trace_transfer}
    An MPO tensor $M$ has zero correlation length when its physical-trace
    transfer is idempotent:
    \begin{align}
        \mathcal T_M^2 & = \mathcal T_M.
        \notag
    \end{align}
    This is the literal fixed-tensor condition of
    \cite[Definition~4.2, lines~735--739]{Cirac2016MPDO_arXiv}.
\end{definition}

\begin{theorem}[Equation for literal physical-trace idempotence]
    \label{thm:mpo_physical_trace_idempotence_iff}
    \lean{MPOTensor.isPhysicalTraceIdempotent_iff}
    \leanok
    \uses{def:mpo_physical_trace_idempotence}
    An MPO tensor satisfies literal physical-trace idempotence if and only if
    $\mathcal T_M^2=\mathcal T_M$.
\end{theorem}

\begin{proof}
    \leanok
    This is the defining equation.
\end{proof}

\begin{definition}[Periodic two-point MPDO correlation]
    \label{def:mpdo_periodic_two_point_correlation}
    \lean{MPOTensor.periodicTwoPointCorrelation}
    \leanok
    \uses{def:mpdo_vertical_loops, def:mpo_physical_trace_transfer}
    For one-site observables $O_1,O_2\in\MN{d}$ and complementary gap
    lengths $g_{\mathrm{mid}},g_{\mathrm{wrap}}\geq0$, set
    \begin{align}
        C_M(O_1,O_2;g_{\mathrm{mid}},g_{\mathrm{wrap}})
         & =\tr\!\left(
        \mathbb E_{O_1}\mathcal T_M^{g_{\mathrm{mid}}}
        \mathbb E_{O_2}\mathcal T_M^{g_{\mathrm{wrap}}}\right),
        \label{eq:mpdo_periodic_correlation}
    \end{align}
    where
    $\mathbb E_O=\sum_{i,j}O_{ji}M^{ij}$.
    In the notation of the two-observable formula in
    \cite[lines~490--496]{Cirac2016MPDO_arXiv},
    $g_{\mathrm{mid}}=n_2-n_1-1$,
    $g_{\mathrm{wrap}}=N+n_1-n_2-1$, and
    $N=g_{\mathrm{mid}}+g_{\mathrm{wrap}}+2$.
    This is the unnormalized contraction: the source introduces normalization
    only at lines~792--793 for entropic quantities.
\end{definition}

\begin{theorem}[Zero correlation length removes positive wrapping-gap dependence]
    \label{thm:mpdo_zcl_periodic_correlation_positive_wrap_gaps_independent}
    \lean{MPOTensor.periodicTwoPointCorrelation_positiveWrapGaps_independent}
    \leanok
    \uses{def:mpo_physical_trace_idempotence,
        def:mpdo_periodic_two_point_correlation}
    Let $M$ be an MPO tensor satisfying the literal zero-correlation-length
    condition. Fix $O_1,O_2$ and a middle gap $g_{\mathrm{mid}}\geq0$.
    For two positive wrapping gaps
    $g_{\mathrm{wrap}},g'_{\mathrm{wrap}}>0$,
    \begin{align}
        C_M(O_1,O_2;g_{\mathrm{mid}},g_{\mathrm{wrap}})
         & =C_M(O_1,O_2;g_{\mathrm{mid}},g'_{\mathrm{wrap}}).
        \label{eq:mpdo_zcl_correlation_positive_wrap_gap_independence}
    \end{align}
    This is independence under changing the ring length at fixed observable
    separation. It includes adjacent marked sites, for which
    $g_{\mathrm{mid}}=0$. In particular, it applies when $M$ generates an
    MPDO, which is the source context of the assertion following
    \cite[Definition~4.2, lines~735--742]{Cirac2016MPDO_arXiv}. The exact
    positive-power boundary is recorded in
    \cite{gap:cpsv16_mpdo_zcl_correlation_length_boundary}.
    No invariance between a zero wrapping gap and a positive one is asserted.
\end{theorem}

\begin{proof}
    \leanok
    Idempotence gives $\mathcal T_M^r=\mathcal T_M$ for every $r>0$.
    Apply this identity to the wrapping power in each side of
    (\ref{eq:mpdo_zcl_correlation_positive_wrap_gap_independence}); the middle
    power is unchanged.
\end{proof}

\begin{corollary}[Zero correlation length gives positive-gap distance independence]
    \label{thm:mpdo_zcl_periodic_correlation_positive_gaps_independent}
    \lean{MPOTensor.periodicTwoPointCorrelation_positiveGaps_independent}
    \leanok
    \uses{def:mpo_physical_trace_idempotence,
        def:mpdo_periodic_two_point_correlation}
    Let $M$ be an MPO tensor satisfying the literal zero-correlation-length
    condition. Fix $O_1,O_2$. For two pairs of positive complementary gaps
    $g_{\mathrm{mid}},g_{\mathrm{wrap}}>0$ and
    $g'_{\mathrm{mid}},g'_{\mathrm{wrap}}>0$,
    \begin{align}
        C_M(O_1,O_2;g_{\mathrm{mid}},g_{\mathrm{wrap}})
         & =C_M(O_1,O_2;g'_{\mathrm{mid}},g'_{\mathrm{wrap}}).
        \label{eq:mpdo_zcl_correlation_positive_gap_independence}
    \end{align}
    For equal gap sums this is correlation independence under changing the
    observable distance on a fixed ring while the observables are nonadjacent
    on both complementary arcs, as in the CID discussion at
    \cite[lines~433--500]{Cirac2016MPDO_arXiv}. The positive-power boundary is
    recorded in
    \cite{gap:cpsv16_mpdo_zcl_correlation_length_boundary}.
    No invariance between a zero complementary gap and a positive one is
    asserted.
\end{corollary}

\begin{proof}
    \leanok
    Idempotence gives $\mathcal T_M^r=\mathcal T_M$ for every $r>0$.
    Applying this identity to both powers in each instance of
    (\ref{eq:mpdo_periodic_correlation}) proves
    (\ref{eq:mpdo_zcl_correlation_positive_gap_independence}).
\end{proof}

\begin{definition}[Nonzero up-to-scalar physical-trace relation]
    \label{def:mpo_source_zcl}
    \lean{MPOTensor.IsSourceZCL}
    \leanok
    \uses{def:mpo_physical_trace_transfer}
    An MPO tensor satisfies the nonzero up-to-scalar physical-trace relation
    when, for some real number $\lambda>0$, the following conditions hold:
    \begin{align}
        \mathcal T_M   & \ne 0,
        \notag                                   \\
        \mathcal T_M^2 & =\lambda\,\mathcal T_M.
        \notag
    \end{align}
    The nonzero condition excludes the degenerate zero transfer, and the
    positive scalar makes the condition invariant under positive real
    rescaling of $M$.
    This is a scale-invariant relation at the level of the generated state
    family. It is broader than the paper's literal diagram, which fixes
    $\lambda=1$ for the displayed tensor.
    % **Scope restriction (up-to-scalar interpretation):** documented in
    % docs/paper-gaps/cpsv16_zcl_canonical_form_normalization.tex.
\end{definition}

\begin{theorem}[The nonzero physical-trace relation requires a nonzero bond
        space]
    \label{thm:mpo_source_zcl_bond_dim_ne_zero}
    \lean{MPOTensor.IsSourceZCL.bondDim_ne_zero}
    \leanok
    \uses{def:mpo_source_zcl}
    If an MPO tensor satisfies the nonzero up-to-scalar physical-trace
    relation, then its bond dimension is positive.
\end{theorem}

\begin{proof}\leanok
    On a zero-dimensional bond space every endomorphism vanishes, contrary to
    the nonzero physical-trace transfer in
    Definition~\ref{def:mpo_source_zcl}.
\end{proof}

\begin{theorem}[Literal idempotence gives the up-to-scalar physical-trace
        relation]
    \label{thm:mpo_source_zcl_of_physical_trace_idempotent}
    \lean{MPOTensor.IsPhysicalTraceIdempotent.isSourceZCL}
    \leanok
    \uses{def:mpo_source_zcl, def:mpo_physical_trace_idempotence}
    If $\mathcal T_M\ne0$ and $M$ has literal physical-trace idempotence, then $M$
    satisfies the nonzero up-to-scalar physical-trace relation.
\end{theorem}

\begin{proof}
    \leanok
    This is the preceding definition with $\lambda=1$.
\end{proof}

\begin{align}
    \mathcal T_M & =\mathcal T_M^2.
    \notag
\end{align}
\begin{center}
    % T_M = T_M^2: the physical-trace transfer T_M = sum_i M^{ii} is
    % idempotent.  Each site is one MPDO tensor M whose ket leg is
    % arced into its own bra leg (trace=physical); the wire between
    % the two sites on the right is the contracted virtual index of
    % the product T_M T_M.  Traced physical legs are closed indices,
    % so both sides carry the same boundary signature
    % (virtual-west=1 | virtual-east=1 | physical-up=0 | physical-down=0).
    \tenkzkernel
    \begin{tenkz}[rows={wire}, cols=1, west=open, east=open, trace=physical,
            physical=updown]
        \tn[skin=mpo]{M}
    \end{tenkz}
    $ = $
    \begin{tenkz}[rows={wire}, cols=2, west=open, east=open, trace=physical,
            physical=updown]
        \tn[skin=mpo]{M} & \tn[skin=mpo]{M}
    \end{tenkz}
\end{center}
This is the literal identity in
\cite[Definition~4.2, lines~735--739]{Cirac2016MPDO_arXiv};
Definition~\ref{def:mpo_source_zcl} permits the positive factor $\lambda$.

\begin{lemma}[Physical-trace normalization gives an idempotent transfer]
    \label{lem:mpo_source_zcl_normalized_idempotent}
    \lean{MPOTensor.IsSourceZCL.normalized_idempotent}
    \leanok
    \uses{def:mpo_source_zcl}
    If $M$ satisfies the nonzero up-to-scalar physical-trace relation, then
    there exists $\lambda>0$ such that
    $\lambda^{-1}\mathcal T_M$ is idempotent.
\end{lemma}

\begin{proof}
    \leanok
    Choose $\lambda>0$ from the up-to-scalar relation, so that
    $\mathcal T_M^2=\lambda\,\mathcal T_M$. Then
    \begin{align}
        (\lambda^{-1}\mathcal T_M)^2
         & =\lambda^{-2}\mathcal T_M^2
        =\lambda^{-1}\mathcal T_M.
        \notag
    \end{align}
\end{proof}

\begin{definition}[Doubled-index transfer idempotence for MPO tensors]%
    \label{def:mpo_is_zcl}
    \lean{MPOTensor.IsZCL}
    \leanok
    \uses{def:mpo_transfer_map}
    An MPO tensor $M$ satisfies the doubled-index transfer condition when its
    completely positive transfer map is idempotent:
    \begin{align}
        \E_M\circ\E_M & =\E_M.
        \notag
    \end{align}
    % Gap documented in docs/paper-gaps/cpsv16_zcl_canonical_form_normalization.tex.
    This is the condition obtained from the doubled-index MPS view. It is not
    the physical-trace condition of
    Definition~\ref{def:mpo_physical_trace_idempotence}.
\end{definition}

\begin{theorem}[Doubled-index transfer idempotence iff MPS RFP]%
    \label{thm:mpo_is_zcl_iff_to_mps_rfp}
    \lean{MPOTensor.isZCL_iff_toMPSTensor_isTransferIdempotent}
    \leanok
    \uses{def:mpo_is_zcl, def:mpo_to_mps, def:mps_transfer_idempotence, lem:mpo_transfer_eq_mps}
    An MPO tensor satisfies the doubled-index transfer condition if and only if
    its doubled-index MPS view is a renormalization fixed point.
\end{theorem}

\begin{proof}
    \leanok
    By Lemma~\ref{lem:mpo_transfer_eq_mps}, the transfer map of $M$ agrees with
    the transfer map of its doubled-index MPS view. Thus
    $\E_M\circ\E_M=\E_M$ holds exactly when the transfer map of
    that doubled-index MPS tensor is idempotent, which is the RFP condition.
\end{proof}

\begin{theorem}[Renormalization fixed points have idempotent physical-trace transfer]
    \label{thm:rfp_via_ts_physical_trace_idempotent}
    \lean{MPOTensor.isPhysicalTraceIdempotent_of_isRFPViaTS}
    \leanok
    \uses{def:is_rfp_via_ts, def:mpo_physical_trace_idempotence}
    If $M$ is a renormalization fixed point via trace-preserving maps, then
    $\mathcal T_M^2=\mathcal T_M$.
    This proves the zero-correlation-length component of implication
    (i)$\Rightarrow$(ii) in
    \cite[Theorem~4.9, lines~851--893]{Cirac2016MPDO_arXiv}. The corresponding
    calculation appears in the proposition ``RFP implies ZCL and SAL'' in
    Appendix~C, lines~1333--1340; the part establishing saturation of the area
    law is separate.
\end{theorem}

\begin{proof}
    \leanok
    For every virtual matrix $X$, the one-site and two-site closures satisfy
    \begin{align}
        \tr M_1(X) & =\tr(\mathcal T_M X),
        \notag                               \\
        \tr M_2(X) & =\tr(\mathcal T_M^2 X).
        \notag
    \end{align}
    Write $\Phi$ for the one-to-two trace-preserving map $\mathcal T$ of
    Definition~\ref{def:is_rfp_via_ts}; this is distinct from the
    physical-trace transfer $\mathcal T_M$. Then
    $\Phi[M_1(X)]=M_2(X)$, and hence
    \begin{align}
        \tr(\mathcal T_M^2 X)
         & =\tr M_2(X)
        =\tr\Phi[M_1(X)]
        =\tr M_1(X)
        =\tr(\mathcal T_M X).
        \notag
    \end{align}
    Nondegeneracy of the trace pairing now gives
    $\mathcal T_M^2=\mathcal T_M$.
\end{proof}

\begin{theorem}[An idempotent physical-trace transfer without fixed-tensor
        renormalization maps]
    \label{thm:kato_deformed_mpdo_fixed_tensor_obstruction}
    \lean{MPOTensor.KatoDeformedRFPObstruction.pauliZ,
        MPOTensor.siteSign,
        MPOTensor.KatoDeformedRFPObstruction.tensor,
        MPOTensor.configurationSign,
        MPOTensor.KatoDeformedRFPObstruction.mpo_tensor_eq_diagonal,
        MPOTensor.KatoDeformedRFPObstruction.tensor_isMPDO,
        MPOTensor.KatoDeformedRFPObstruction.trace_mpo_tensor,
        MPOTensor.KatoDeformedRFPObstruction.physTraceTransfer_tensor,
        MPOTensor.KatoDeformedRFPObstruction.%
        physTraceTransfer_tensor_idempotent,
        MPOTensor.KatoDeformedRFPObstruction.physClose1_tensor,
        MPOTensor.KatoDeformedRFPObstruction.physClose2_tensor,
        MPOTensor.KatoDeformedRFPObstruction.obstructionBoundary,
        MPOTensor.KatoDeformedRFPObstruction.%
        physClose1_obstructionBoundary,
        MPOTensor.KatoDeformedRFPObstruction.%
        physClose2_obstructionBoundary,
        MPOTensor.KatoDeformedRFPObstruction.%
        physClose2_obstructionBoundary_posSemidef,
        MPOTensor.KatoDeformedRFPObstruction.%
        physClose1_obstructionBoundary_not_posSemidef,
        MPOTensor.KatoDeformedRFPObstruction.tensor_not_isRFPViaTS}
    \leanok
    \uses{def:mpdo, def:phys_close, def:mpo_physical_trace_transfer, def:is_rfp_via_ts}
    Let $Z=\operatorname{diag}(1,-1)$ and let $K$ be the bond-two MPO tensor
    whose nonzero physical letters are
    \begin{align}
        K^{00} & =\operatorname{diag}(1/2,1/4),  &
        K^{11} & =\operatorname{diag}(1/2,-1/4).
        \notag
    \end{align}
    For every $N\geq1$, its closed operator is
    \begin{align}
        \rho_N(K)    & =2^{-N}I+4^{-N}Z^{\otimes N}\geq0, &
        \tr\rho_N(K) & =1.
        \notag
    \end{align}
    Its physical-trace transfer is
    \begin{align}
        \mathcal T_K   & =\operatorname{diag}(1,0), &
        \mathcal T_K^2 & =\mathcal T_K.
        \notag
    \end{align}
    For every virtual matrix $X$, the physical closures satisfy
    \begin{align}
        K_1(X) & =\frac{X_{00}}2I+\frac{X_{11}}4Z,
        \notag                                     \\
        K_2(X) & =\frac{X_{00}}4I\otimes I
        +\frac{X_{11}}{16}Z\otimes Z.
        \notag
    \end{align}
    Nevertheless, $K$ is not a renormalization fixed point in the sense of
    Definition~\ref{def:is_rfp_via_ts}.

    The tensor is the $p=1/2$ example of
    \cite[lines~709--838]{Kato2024Exact}, after conjugating every physical site
    by the Hadamard matrix. Kato's exact renormalization changes the tensor
    along the flow and the source explicitly states that this example is not a
    fixed point. The fixed-tensor-channel obstruction follows directly from
    the one- and two-site closure identities; it is not asserted as a theorem
    in that source. No canonical-form or strong-area-law conclusion is part of
    this statement. Theorem~\ref{thm:kato_deformed_mpdo_cf_sal_capstone}
    establishes both conclusions for the same tensor.
\end{theorem}

\begin{proof}
    \leanok
    \uses{lem:mpdo_trace_loop_chain}
    Since the two virtual sectors are diagonal, set $z_0=1$ and $z_1=-1$.
    A physical configuration $\sigma$ has weight
    \begin{align}
        2^{-N}+4^{-N}\prod_{k=0}^{N-1}z_{\sigma_k}.
        \notag
    \end{align}
    This proves the closed-operator formula. Each diagonal entry is either
    $2^{-N}+4^{-N}$ or $2^{-N}-4^{-N}$ and is nonnegative. Contracting the
    physical indices gives $\mathcal T_K=\operatorname{diag}(1,0)$, whence
    $\tr\rho_N(K)=\tr(\mathcal T_K^N)=1$ for $N\geq1$ and
    $\mathcal T_K^2=\mathcal T_K$.

    Direct contraction with a virtual matrix $X$ gives the two closure
    formulas. Set $X_* = \operatorname{diag}(2,8)$. Then
    \begin{align}
        K_2(X_*) & =\frac12(I\otimes I+Z\otimes Z)\geq0,
        \notag                                               \\
        K_1(X_*) & =I+2Z=\operatorname{diag}(3,-1)\not\geq0.
        \notag
    \end{align}
    If the coarse-graining channel $\mathcal S$ of
    Definition~\ref{def:is_rfp_via_ts} existed, complete positivity would give
    $\mathcal S[K_2(X_*)]\geq0$, whereas its intertwining identity would give
    $\mathcal S[K_2(X_*)]=K_1(X_*)$. This is a contradiction.
\end{proof}

\begin{theorem}[A canonical-form, area-law-saturating MPDO without
        fixed-tensor renormalization maps]
    \label{thm:kato_deformed_mpdo_cf_sal_capstone}
    \lean{MPOTensor.KatoDeformedRFPObstruction.%
        tensor_toMPSTensor_isCPSVCanonicalForm,
        MPOTensor.KatoDeformedRFPObstruction.tensor_isSAL,
        MPOTensor.KatoDeformedRFPObstruction.%
        exists_isCPSVCanonicalForm_isMPDO_idempotent_isSAL_not_isRFPViaTS}
    \leanok
    \uses{thm:kato_deformed_mpdo_fixed_tensor_obstruction, def:mpdo, def:mpo_physical_trace_transfer, def:is_rfp_via_ts, def:is_sal, def:cpsv_canonical_form}
    Let $K$ be the bond-two MPO tensor of
    Theorem~\ref{thm:kato_deformed_mpdo_fixed_tensor_obstruction}, and let
    $A_K$ be its doubled-index MPS tensor, an MPS tensor with four physical
    symbols. Then $A_K$ is in literal CPSV canonical form: every letter is
    the weighted direct sum
    $A_K^{(i,j)}=\mu_0A_0^{(i,j)}\oplus\mu_1A_1^{(i,j)}$ for
    $i,j\in\{0,1\}$, with weights and bond-one normal tensors
    \begin{align}
        \mu_0       & =\frac1{\sqrt2},                     &
        \mu_1       & =\frac1{2\sqrt2},
        \notag                                               \\
        A_0^{(i,j)} & =\frac{\delta_{ij}}{\sqrt2},         &
        A_1^{(i,j)} & =(-1)^i\,\frac{\delta_{ij}}{\sqrt2}.
        \notag
    \end{align}
    Moreover, $K$ saturates the area law: for every $N\geq2$ and every
    $1\leq L<N$, the reduced state of the first $L$ sites is maximally
    mixed,
    \begin{align}
        \rho_L^{(N)}(K) & =2^{-L}I, &
        S_L^{(N)}(K)    & =L\log2.
        \notag
    \end{align}
    Hence the mutual information
    $I_L^{(N)}=S_L^{(N)}(K)+S_{N-L}^{(N)}(K)-S_N^{(N)}(K)
        =N\log2-S_N^{(N)}(K)$ is constant in $L$, giving
    $I_1=I_2=\cdots=I_{\lfloor N/2\rfloor}$.
    Consequently, there is a bond-two MPO tensor whose doubled-index MPS
    tensor is in literal CPSV canonical form, which generates an MPDO, whose
    physical-trace transfer is idempotent, which saturates the area law, and
    which is not a renormalization fixed point in the sense of
    Definition~\ref{def:is_rfp_via_ts}.

    The canonical-form decomposition follows the construction of
    \cite[Section~2.3, lines~214--245]{Cirac2016MPDO_arXiv}. The
    decomposition and the entropy evaluation are project derivations;
    neither is asserted for this tensor in \cite{Kato2024Exact} or in
    \cite{Cirac2016MPDO_arXiv}.
\end{theorem}

\begin{proof}
    \leanok
    \uses{thm:kato_deformed_mpdo_fixed_tensor_obstruction, def:cpsv_canonical_form, def:is_sal, lem:bond_dimension_one_identity_transfer_normal}
    Every letter of $A_K$ is diagonal and the block letters are scalars, so
    the weighted direct sum is the entrywise check
    \begin{align}
        \mu_0A_0^{(i,i)} & =\frac12,        &
        \mu_1A_1^{(i,i)} & =\frac{(-1)^i}4,
        \notag
    \end{align}
    whose right-hand sides are the two virtual diagonal entries of $K^{ii}$;
    the letters with $i\ne j$ vanish on both sides. Each block satisfies
    $\sum_{i,j}A_k^{(i,j)}\,\overline{A_k^{(i,j)}}=1$, hence has identity
    transfer map and is a normal tensor by
    Lemma~\ref{lem:bond_dimension_one_identity_transfer_normal}. The
    weighted direct sum of positive-dimensional normal blocks is in literal
    CPSV canonical form by Definition~\ref{def:cpsv_canonical_form}.

    Since $\tr\rho_N(K)=1$, the normalized state is $\rho_N(K)$ itself.
    For $1\leq L<N$, the partial trace over the complementary $N-L$ sites
    annihilates the summand $4^{-N}Z^{\otimes N}$ of
    $\rho_N(K)=2^{-N}I+4^{-N}Z^{\otimes N}$, since $\tr Z=0$, and
    contracts the remaining summand to $2^{-L}I$. Hence
    $\rho_L^{(N)}(K)=2^{-L}I$ and $S_L^{(N)}(K)=L\log2$. Substitution in
    $I_L^{(N)}=S_L^{(N)}(K)+S_{N-L}^{(N)}(K)-S_N^{(N)}(K)$ gives
    $I_L^{(N)}=N\log2-S_N^{(N)}(K)$, which is independent of
    $L$; in particular $I_L=I_{L+1}$ whenever
    $1\leq L<\lfloor N/2\rfloor$, so $K$ saturates the area law by
    Definition~\ref{def:is_sal}.

    Theorem~\ref{thm:kato_deformed_mpdo_fixed_tensor_obstruction} supplies
    the remaining three properties: $K$ generates an MPDO, its
    physical-trace transfer is idempotent, and it is not a renormalization
    fixed point. Together with the two properties above, this gives the
    existential.
\end{proof}

\begin{theorem}[Renormalization fixed points satisfy the up-to-scalar
        physical-trace relation]
    \label{thm:rfp_via_ts_source_zcl}
    \lean{MPOTensor.isSourceZCL_of_isRFPViaTS}
    \leanok
    \uses{thm:rfp_via_ts_physical_trace_idempotent, thm:mpo_source_zcl_of_physical_trace_idempotent}
    If $M$ is a renormalization fixed point via trace-preserving maps and
    $\mathcal T_M\ne0$, then $M$ satisfies the nonzero up-to-scalar
    physical-trace relation.
\end{theorem}

\begin{proof}
    \leanok
    Apply the preceding idempotence theorem and take the positive scalar in
    Definition~\ref{def:mpo_source_zcl} to be $1$.
\end{proof}

\section{Purification fixed points}

\begin{definition}[Purifying spin--ancilla tensor]
    \label{def:mpdo_purification_tensor}
    \lean{MPOTensor.purificationTensor}
    \leanok
    \uses{def:mps_tensor}
    Given matrices $A^{(i,k)}$ with a spin index $i$ and an ancillary index $k$,
    the associated MPS tensor on the product physical space is
    $\widehat A^{(i,k)}=A^{(i,k)}$.
    This is the tensor whose finite-chain vector appears in the purification
    formula of \cite[lines~744--751]{Cirac2016MPDO_arXiv}.
\end{definition}

\begin{definition}[Ancillary trace of a purifying MPS]
    \label{def:mpdo_purification_density}
    \lean{MPOTensor.purificationDensity}
    \leanok
    \uses{def:mpdo_purification_tensor, def:mpv}
    For a purifying spin--ancilla tensor $A$, the reduced spin density matrix is
    defined by
    \begin{align}
        [\rho^{(N)}_p(A)]_{\sigma,\tau}
         & =
        \sum_{\kappa}
        \psi_A(\sigma,\kappa)
        \overline{\psi_A(\tau,\kappa)},
        \notag \\
        \psi_A(\sigma,\kappa)
         & =
        \tr\!\left(A^{(\sigma_0,\kappa_0)}
        A^{(\sigma_1,\kappa_1)}
        \cdots
        A^{(\sigma_{N-1},\kappa_{N-1})}\right).
        \notag
    \end{align}
    This is the coefficient form of the ancillary trace in
    \cite[eq.~\textup{(4.7)}, line~751]{Cirac2016MPDO_arXiv}.
\end{definition}

\begin{definition}[Single-site ancillary trace map]
    \label{def:mpdo_ancillary_trace_map}
    \lean{MPOTensor.ancillaryTraceMap}
    \leanok

    For a matrix $X$ on the spin--ancilla space, the ancillary trace map is
    \begin{align}
        [\tr_a(X)]_{ij}
         & =
        \sum_k X_{(i,k),(j,k)}.
        \notag
    \end{align}
    This is the one-site trace over ancillary degrees of freedom appearing in
    \cite[lines~751 and~761--764]{Cirac2016MPDO_arXiv}.
\end{definition}

\begin{theorem}[The ancillary trace is trace-preserving completely positive]
    \label{thm:mpdo_ancillary_trace_map_is_kraus_cptp}
    \lean{MPOTensor.ancillaryTraceMap_isKrausCPTP}
    \leanok
    \uses{def:mpdo_ancillary_trace_map}
    For every ancillary dimension, the one-site ancillary trace map is
    trace-preserving and completely positive.
\end{theorem}
\begin{proof}
    \leanok

    This is Lemma~\cite[``The right partial trace is trace-preserving completely positive'']{QICLean2026Blueprint} with $A$ the spin
    space and $B$ the ancillary space.
\end{proof}

\begin{definition}[Global purification equation]
    \label{def:mpdo_global_purification_equation}
    \lean{MPOTensor.HasGlobalPurificationEquation}
    \leanok
    \uses{def:mpo_family, def:mpdo_purification_density}
    An MPO tensor $M$ satisfies the global purification equation with a
    purifying tensor $A$ if
    $\rho^{(N)}(M)=\rho^{(N)}_p(A)$ for every positive length $N$.
    Equivalently, for every positive length $N$,
    \begin{align}
        \rho^{(N)}(M)
         & =
        \tr_a
        \bigl(        \ketbra{\Psi^{(N)}(A)}{\Psi^{(N)}(A)}
        \bigr),
        \notag
    \end{align}
    as in \cite[line~751]{Cirac2016MPDO_arXiv}.
\end{definition}

\begin{definition}[Global-family purification RFP witness]
    \label{def:mpdo_purification_rfp_witness}
    \lean{MPOTensor.HasPurificationRFPWitness}
    \leanok
    \uses{def:mpdo_global_purification_equation, def:mps_transfer_idempotence}
    An MPO tensor $M$ has a global-family purification RFP witness if, for some
    purifying spin--ancilla tensor $A$, it satisfies the positive-length global
    purification equation and $\widehat A$ is a pure-state renormalization fixed
    point. This auxiliary condition records the density-operator equation at
    \cite[line~751]{Cirac2016MPDO_arXiv}. It does not assert the one-site
    ancillary-contraction presentation in Definition~\ref{def:mpdo_is_prfp}.
\end{definition}

\begin{definition}[Trace-preserving spin reduction]
    \label{def:mpdo_has_trace_preserving_spin_reduction}
    \lean{MPOTensor.HasTracePreservingSpinReduction}
    \leanok
    \uses{def:mpdo_ancillary_trace_map}
    The spin reduction obtained by tracing the ancillary index is
    trace-preserving completely positive: the ancillary trace map
    $\tr_a$ satisfies the corresponding definition in the companion quantum-channel volume~\cite[``Trace-preserving completely positive map'']{QICLean2026Blueprint}.
    This is the structure described immediately after Definition~4.3 in
    \cite[lines~761--764]{Cirac2016MPDO_arXiv}.
\end{definition}

\begin{corollary}[Trace-preserving spin reduction of the ancillary trace]
    \label{cor:mpdo_has_trace_preserving_spin_reduction}
    \lean{MPOTensor.hasTracePreservingSpinReduction}
    \leanok
    \uses{def:mpdo_has_trace_preserving_spin_reduction, thm:mpdo_ancillary_trace_map_is_kraus_cptp}
    For every ancillary dimension, the spin reduction obtained by tracing the
    ancillary index is trace-preserving and completely positive.
\end{corollary}
\begin{proof}
    \leanok
    This is exactly
    Theorem~\ref{thm:mpdo_ancillary_trace_map_is_kraus_cptp}, unpacked through
    Definition~\ref{def:mpdo_has_trace_preserving_spin_reduction}.
\end{proof}

% **Local fix (intended presentation):** Definition 4.3 literally refers to
% the global equation MPDO-Puri-1, while lines 744--763 continue with the
% one-site presentation Psipuri. We adopt the latter standing presentation.
% See docs/paper-gaps/cpsv16_purification_rfp_definition.tex.
\begin{definition}[Purification renormalization fixed point]
    \label{def:mpdo_is_prfp}
    \lean{MPOTensor.IsPRFP}
    \leanok
    \uses{def:lpdo, def:mps_transfer_idempotence}
    An MPO tensor $M$ is a purification renormalization fixed point if there is
    a spin--ancilla tensor $A$ and an identification of the doubled purifying
    bond space with the bond space of $M$ such that
    \begin{align}
        M^{ij}
         & =
        \left(\sum_k A^{(i,k)}\otimes
        \overline{A^{(j,k)}}\right)_{e,e},
        \notag
    \end{align}
    and $\widehat A$ is a pure-state renormalization fixed point. Here PRFP
    denotes the standing one-site presentation; the positive-length equation is
    its MPDO interpretation.
\end{definition}

\begin{definition}[Nondegenerate purification RFP]
    \label{def:is_nondegenerate_prfp}
    \lean{MPOTensor.IsNondegeneratePRFP}
    \leanok
    \uses{def:mpdo_is_prfp, def:mpo_physical_trace_transfer}
    An MPO tensor is a nondegenerate purification renormalization fixed point if
    it is a purification renormalization fixed point and
    $\mathcal T_M=\sum_i M^{ii}\ne0$.
\end{definition}

\begin{theorem}[A purification RFP is an LPDO]
    \label{thm:prfp_is_lpdo}
    \lean{MPOTensor.IsPRFP.isLPDO}
    \leanok
    \uses{def:mpdo_is_prfp, def:lpdo}
    Every purification renormalization fixed point is a local purification
    density operator.
\end{theorem}

\begin{proof}\leanok
    Forget the renormalization-fixed-point condition on the purifying tensor.
\end{proof}

\begin{theorem}[A purification RFP generates MPDOs]
    \label{thm:prfp_is_mpdo}
    \lean{MPOTensor.IsPRFP.isMPDO}
    \leanok
    \uses{def:mpdo_is_prfp, def:mpdo}
    Every purification renormalization fixed point generates matrix product
    density operators.
\end{theorem}

\begin{proof}\leanok
    \uses{thm:prfp_is_lpdo, thm:lpdo_is_mpdo}
    Apply Theorem~\ref{thm:lpdo_is_mpdo} to the local purification supplied by
    Theorem~\ref{thm:prfp_is_lpdo}.
\end{proof}

\begin{theorem}[Purification RFP does not imply doubled-index idempotence]
    \label{thm:prfp_not_doubled_index_zcl}
    \lean{MPOTensor.exists_isPRFP_not_isZCL}
    \leanok
    \uses{def:mpdo_is_prfp, def:mpo_is_zcl}
    There is a purification renormalization fixed point for which the
    doubled-index transfer-map equation $\E_M\circ\E_M=\E_M$ fails.
\end{theorem}

% Counterexample analysed in docs/paper-gaps/cpsv16_zcl_canonical_form_normalization.tex.
\begin{proof}\leanok
    Take the diagonal purification at $d=d_K=2$ and $D=D'=1$ with amplitudes
    $A=[\tfrac{1}{\sqrt2},0,0,\tfrac{1}{\sqrt2}]$. The purifying tensor is a
    pure-state renormalization fixed point, since $\sum|A|^2=1$ makes its transfer
    map the identity. The ancilla contraction
    $M^{ij}=\sum_k A^{(i,k)} \overline{A^{(j,k)}}$
    gives the scalar entries $M^{00}=M^{11}=\tfrac12$ (off-diagonal $0$), so the
    induced transfer map is $\tfrac12\cdot\id$ and
    $\E_M\circ\E_M=\tfrac14\cdot\id\neq\E_M$. The trace contraction has
    dropped the leading eigenvalue from $1$ to $\tfrac12$. By contrast, the
    physical-trace transfer is $1$ and hence is literally idempotent. Thus this
    example separates doubled-index idempotence from the paper's
    physical-trace condition.
\end{proof}

\begin{lemma}[Physical-trace transfer of the maximally mixed witness]
    \label{lem:phys_trace_transfer_witnessM}
    \lean{MPOTensor.physTraceTransfer_witnessM}
    \leanok
    \uses{def:mpo_physical_trace_transfer}
    For the maximally mixed witness $M_{\mathrm w}$, closing the ket and bra
    physical legs gives
    $\mathcal T_{M_{\mathrm w}}
        =M_{\mathrm w}^{00}+M_{\mathrm w}^{11}
        =\tfrac12+\tfrac12=1$.
\end{lemma}

\begin{proof}
    \leanok
    This is the direct computation of the two nonzero diagonal entries of the
    witness tensor.
\end{proof}

\begin{theorem}[The maximally mixed witness satisfies the nonzero
        physical-trace relation]
    \label{thm:witness_source_zcl}
    \lean{MPOTensor.isSourceZCL_witnessM}
    \leanok
    \uses{def:mpo_source_zcl, thm:prfp_not_doubled_index_zcl}
    The maximally mixed witness $M_{\mathrm w}$ satisfies the nonzero
    up-to-scalar physical-trace relation.
\end{theorem}
\begin{proof}\leanok
    \uses{lem:phys_trace_transfer_witnessM, thm:mpo_source_zcl_of_physical_trace_idempotent}
    Lemma~\ref{lem:phys_trace_transfer_witnessM} gives
    $\mathcal T_{M_{\mathrm w}}=1$. The identity is nonzero and idempotent, so
    the relation holds with $\lambda=1$.
\end{proof}

\begin{theorem}[Fixed local purification: pure RFP iff normalized physical-trace idempotence]
    \label{thm:purification_tensor_rfp_iff_physical_trace_idempotent}
    \lean{MPOTensor.purificationTensor_isTransferIdempotent_iff_physTraceTransfer_sq}
    \leanok
    \uses{def:mpdo_purification_tensor, def:mpo_physical_trace_transfer, def:mps_transfer_idempotence}
    Suppose that $M$ is the ancillary contraction of a fixed spin--ancilla
    tensor $A$ through a bond identification $e$:
    \begin{align}
        M^{ij}
         & =
        \left(\sum_k A^{(i,k)}\otimes
        \overline{A^{(j,k)}}\right)_{e,e}.
        \notag
    \end{align}
    Then the purification tensor $\widehat A$ is a pure-state renormalization
    fixed point if and only if $\mathcal T_M^2=\mathcal T_M$.
    This is the tensor-level equivalence between the first two conditions of
    the purification theorem in \cite[lines~744--786]{Cirac2016MPDO_arXiv}.
\end{theorem}

\begin{proof}
    \leanok
    Let $K'$ be the transfer matrix of $\widehat A$, let $s$ interchange the
    two Kronecker factors, and let $K$ denote the same Kronecker sum in the
    physical ordering. The local contraction gives
    \begin{align}
        \mathcal T_M & =K_e,
        \notag                  \\
        K            & =(K')_s.
        \notag
    \end{align}
    Since both reindexings are bijective and commute with matrix
    multiplication,
    \begin{align}
        (K')^2=K'
        \quad & \Longleftrightarrow\quad
        \mathcal T_M^2=\mathcal T_M.
        \notag
    \end{align}
    Injectivity of the transfer-matrix representation gives
    \begin{align}
        (K')^2=K'
        \quad & \Longleftrightarrow\quad
        \E_{\widehat A}\circ\E_{\widehat A}
        =\E_{\widehat A}.
        \notag
    \end{align}
\end{proof}

\begin{theorem}[A purification RFP has idempotent physical-trace transfer]
    \label{thm:prfp_physical_trace_idempotent}
    \lean{MPOTensor.physTraceTransfer_sq_of_isPRFP}
    \lean{MPOTensor.IsPRFP.isPhysicalTraceIdempotent}
    \leanok
    \uses{def:mpdo_is_prfp, def:mpo_physical_trace_idempotence}
    If $M$ is a purification renormalization fixed point, then
    \begin{align}
        \mathcal T_M^2 & = \mathcal T_M.
        \notag
    \end{align}
    Thus the forward implication from PRFP to the literal zero-correlation-
    length equation in \cite[Theorem~4.4]{Cirac2016MPDO_arXiv} requires no
    nonzero assumption.
\end{theorem}

\begin{proof}\leanok
    \uses{thm:purification_tensor_rfp_iff_physical_trace_idempotent}
    Apply the forward implication of
    Theorem~\ref{thm:purification_tensor_rfp_iff_physical_trace_idempotent} to
    the purifying tensor and bond identification in the PRFP presentation.
\end{proof}

\begin{theorem}[Purification form and physical-trace idempotence characterize PRFP]
    \label{thm:prfp_iff_lpdo_physical_trace_idempotent}
    \lean{MPOTensor.isPRFP_iff_isLPDO_and_physTraceTransfer_sq}
    \leanok
    \uses{def:mpdo_is_prfp, def:lpdo, def:mpo_physical_trace_idempotence}
    An MPO tensor is a purification renormalization fixed point if and only if
    it has a local purification and its physical-trace transfer is idempotent:
    \begin{align}
        \operatorname{PRFP}(M)
        \quad & \Longleftrightarrow\quad
        \operatorname{LPDO}(M)\ \text{and}\ \mathcal T_M^2=\mathcal T_M.
        \notag
    \end{align}
\end{theorem}

\begin{proof}
    \leanok
    \uses{thm:prfp_is_lpdo, thm:prfp_physical_trace_idempotent,
        thm:purification_tensor_rfp_iff_physical_trace_idempotent}
    The forward direction combines Theorems~\ref{thm:prfp_is_lpdo} and
    \ref{thm:prfp_physical_trace_idempotent}. Conversely, choose the local
    purifying tensor and bond identification and apply
    Theorem~\ref{thm:purification_tensor_rfp_iff_physical_trace_idempotent} to
    that fixed purification.
\end{proof}

\begin{theorem}[Nondegenerate PRFP equivalence]
    \label{thm:nondegenerate_prfp_iff_lpdo_normalized_zcl}
    \lean{MPOTensor.isNondegeneratePRFP_iff}
    \leanok
    \uses{def:is_nondegenerate_prfp, def:lpdo, def:mpo_physical_trace_idempotence}
    An MPO tensor is a nondegenerate purification renormalization fixed point if
    and only if it is a local purification density operator and
    \begin{align}
        \mathcal T_M   & \ne0,
        \notag                          \\
        \mathcal T_M^2 & =\mathcal T_M.
        \notag
    \end{align}
\end{theorem}

\begin{proof}
    \leanok
    \uses{thm:prfp_iff_lpdo_physical_trace_idempotent}
    Combine Theorem~\ref{thm:prfp_iff_lpdo_physical_trace_idempotent} with the
    nonzero clause in Definition~\ref{def:is_nondegenerate_prfp}.
\end{proof}

\begin{theorem}[A nonzero purification RFP satisfies the up-to-scalar relation]
    \label{thm:source_zcl_of_prfp}
    \lean{MPOTensor.isSourceZCL_of_isPRFP}
    \leanok
    \uses{def:mpdo_is_prfp, def:mpo_source_zcl}
    If $M$ is a purification renormalization fixed point and
    $\mathcal T_M\ne0$, then $M$ satisfies the nonzero up-to-scalar
    physical-trace relation.
\end{theorem}

\begin{proof}\leanok
    \uses{thm:prfp_physical_trace_idempotent, thm:mpo_source_zcl_of_physical_trace_idempotent}
    Theorem~\ref{thm:prfp_physical_trace_idempotent} gives literal
    idempotence, and the nonzero hypothesis gives the scalar-one case of
    Definition~\ref{def:mpo_source_zcl}.
\end{proof}

\begin{theorem}[A nondegenerate purification RFP satisfies the nonzero
        physical-trace relation]
    \label{thm:nondegenerate_prfp_source_zcl}
    \lean{MPOTensor.IsNondegeneratePRFP.isSourceZCL}
    \leanok
    \uses{def:is_nondegenerate_prfp, def:mpo_source_zcl}
    Every nondegenerate purification renormalization fixed point satisfies the
    nonzero up-to-scalar physical-trace relation.
\end{theorem}

\begin{proof}
    \leanok
    \uses{thm:source_zcl_of_prfp}
    Apply Theorem~\ref{thm:source_zcl_of_prfp} to the two clauses of
    Definition~\ref{def:is_nondegenerate_prfp}.
\end{proof}

\begin{theorem}[The tensor purification identity yields the global purification equation]
    \label{thm:mpo_eq_purification_density}
    \lean{MPOTensor.mpo_eq_purificationDensity}
    \leanok
    \uses{def:lpdo, def:mpo_family, def:mpdo_purification_density}
    Let $M$ be the ancilla contraction of a purifying spin--ancilla tensor $A$
    through a bond identification $e$, so that
    \begin{align}
        M^{ij}
         & =
        \left(\sum_{k} A^{(i,k)}\otimes
        \overline{A^{(j,k)}}\right)_{e,e}.
        \notag
    \end{align}
    Then at every length $N$ the matrix product operator equals the ancillary
    trace of the pure spin--ancilla matrix product state,
    $\rho^{(N)}(M)=\rho^{(N)}_p(A)$.
    This is the coefficient form of the purification equation of
    \cite[line~751]{Cirac2016MPDO_arXiv}.
\end{theorem}

\begin{proof}\leanok
    Applying the mixed-product property
    $(A\otimes B)(C\otimes D)=(AC)\otimes(BD)$ repeatedly, the length-$N$
    product of the contracted tensor entries distributes over the ancillary
    index sum:
    \begin{align}
        \prod_{l=0}^{N-1}
        \left(\sum_{k}
        A^{(\sigma_l,k)}
        \otimes
        \overline{A^{(\tau_l,k)}}\right)
         & =
        \sum_{\kappa\colon[N]\to[d_K]}
        \left(\prod_l A^{(\sigma_l,\kappa_l)}\right)
        \otimes
        \overline{
            \left(\prod_l A^{(\tau_l,\kappa_l)}\right)
        }.
        \notag
    \end{align}
    Tracing this sum and using
    $\tr(X\otimes\overline Y)=\tr(X) \overline{\tr(Y)}$ gives the coefficient
    of the ancillary trace at every $\sigma,\tau$. The equation holds at every
    length, so no positive-length restriction is needed.
\end{proof}

\begin{theorem}[A purification RFP supplies a global-family witness]
    \label{thm:prfp_has_purification_rfp_witness}
    \lean{MPOTensor.IsPRFP.hasPurificationRFPWitness}
    \leanok
    \uses{def:mpdo_is_prfp, def:mpdo_purification_rfp_witness}
    Every purification renormalization fixed point has a global-family
    purification RFP witness with the same purifying tensor.
\end{theorem}

\begin{proof}\leanok
    \uses{thm:mpo_eq_purification_density}
    The local ancillary-contraction identity implies equality of the two
    generated density operators at every length by
    Theorem~\ref{thm:mpo_eq_purification_density}.
\end{proof}

% CPSV16 Theorem 4.4, the unlabeled theorem environment at lines 777--784.
% Gap: docs/paper-gaps/cpsv16_purification_rfp_definition.tex.
\begin{theorem}[Purification fixed-point equivalence]
    \label{thm:cpsv_theorem44_printed_status}
    \uses{def:mpdo_is_prfp, def:mpo_physical_trace_idempotence}
    \notready
    For an MPO tensor with the one-site purification presentation displayed in
    \cite[lines~744--763]{Cirac2016MPDO_arXiv}, Theorem~4.4 asserts the
    equivalence of the following conditions: the purifying tensor is a
    pure-state renormalization fixed point, the physical-trace transfer is
    idempotent, and the generated density operators have the repeated-copy
    form displayed in \cite[lines~765--773]{Cirac2016MPDO_arXiv}.

    % The identification with the repeated-copy density formula is the only
    % part not yet formalized.
    The equivalence between the first two conditions is
    Theorem~\ref{thm:prfp_iff_lpdo_physical_trace_idempotent}. The
    positive-length global purification equation follows from the local
    presentation by
    Theorem~\ref{thm:prfp_has_purification_rfp_witness}, but it does not by
    itself determine that presentation.
\end{theorem}
